\documentclass[5pt, a4paper]{article}

\usepackage[utf8]{inputenc}
\usepackage[T1]{fontenc}
\usepackage{textcomp}
\usepackage{amsmath, amssymb}
\usepackage{multicol}
\usepackage{proof}
\usepackage{enumitem}
\usepackage{fancyhdr}
\usepackage[margin=1in]{geometry}
% figure support
\usepackage{import}
\usepackage{xifthen}
\pdfminorversion=7

\usepackage[most]{tcolorbox}
\usepackage{pdfpages}
\usepackage{transparent}
\newcommand{\incfig}[1]{%
    \def\svgwidth{\columnwidth}
    \import{./figures/}{#1.pdf_tex}
}

\tcbset{
colback=gray!5,
colframe=blue!50!white,
coltitle=black,
arc=0pt,
boxrule=-0.1mm,
fonttitle=\bfseries,
coltitle=black,
left=2.5mm,
leftrule=1mm,
right=3.5mm,
colbacktitle=black!10!white,
toptitle=0.75mm,
bottomtitle=0.5mm,
}
\begin{document}
\section{Computer Arithmetics}
This section concerns itself mostly with how a computer does arithmetics and
the considerations of a computation problem. We'll cover the concept of
precision and accuracy first. Consider this: 

\begin{tcolorbox}[title=Precision vs Accuracy]
    \textit{Accuracy} of the output of the calculation is the measure of how far
    it is from the true value. The \textit{precision} of a calculation refers to
    how many significant digits (significand) that the output has, regardless
    of it's accuracy.
\end{tcolorbox}

Let's talk a little bit about the problems we are trying to solve. A problem is
refered to as well-posed if a solution exist, unique and varies continuously
with the problem data. Do note that this property is seperate from our algorithm
and although we would like our problems to be well-posed, some problems aren't.

Let's talk a bit more on the computation itself. In general, we would like a
computation to not deviate far from the true value. Of course, having it be zero
is largely unrealistic in parctice, since in order fo that to happen, we would
like to already know the true value, and as such have no need for calculation.

This deviation from the true value, which we would call the error, largely
derive from two main sources: The error that arises from the data itself, and
the error that arises from our computation. 

Take as an example, the computation that one would do when calculating the
surface area of earth. To simplify the computation, we model the earth as a
perfect sphere, and as such \[
    A = 4\pi r^2
\] 
applies. Suppose that we are given the value of $\hat{r}$ as the radius of the
earth. Then, we might arive a some value value $\hat{A}$. However, this value
might deviate from the true value for a number of reasons: 
\begin{enumerate}
    \item Inaccuracy in the computation due to rounding and truncation.
    \item Inaccuracy in the given data ($\hat{r}$)
\end{enumerate}
The first one is cinsidered computational error, while the second belongs to data
error. Since we are more interested in the factor we can control, we would be
focusing more on the first one.

When talking about error, we refer to the absolute difference between our
solution and the true value as the forward error. Similarly, we refer tio the
absolute difference in the input is refered to as a backward error.

Generally, when designing an algorithm to solve a problem, we would like our
problems to be \textit{stable}. As in, small pertrubation in the input would
likewise cause small variance in the output. We quantify this using the
condition number \[
    \kappa = \left| \frac{\text{forward error}}{\text{backward error}}\right|
.\] 
We can also say that \[
| \text{forward error} | = \kappa \times | \text{backward error} |
.\] 
As in, the condition number is the multiplier of error. We generally wan this
number to be as low as possible, though in most cases, the minimum is one. 

Do note that the condition number is a property  of your algorithm, not hte
problem itself. A problem is said to be well conditioned if the condition number
is small. Similarly a problem is called ill-conditioned if the condition number
is large.

\subsection*{Floating Point Arithmetics}

At it's core, computer is a bunch of ones and zeros. To oversimplify a bit, the
process involved in computations requires a way to represents numbers in a
format that is precise, compact and minimizes errors. The standard way of doing
so is through a floating point system. Consider the following sequence of
characters in 1's and 0's: \[
    0.00000000.00000000000000000000000
.\] 
The above is a sequence of numbers in base-2 of length 32. A (single) precision
floating point system represents number in the above format. While there are
other format, this chapter would only cover this format, and also the half
precision floating point, which consists of 16 digits instead of 32.

The first digit, represents the sign of the number, with 0 being positive and 1
being negative. The second set of digits, which consists of 8 digits, determines
the exponent of the digit. The third represents the significand of the number.
\end{document}
